%=============================================================================!
% 13.2  RESCALING AND WEAK COUPLING                                           !
%=============================================================================!
\section{Rescaling and weak coupling}
\label{sec:th-berendsen}

The simplest way to hold a temperature is to multiply every velocity by
the factor that brings the kinetic energy to its target. This section
builds that, and the gentler version of Berendsen and his colleagues,
measures what both do to the distribution of $K$, and derives two faults
they share.

\subsection*{Rescaling}

Before each step, \texttt{Rescale} multiplies every velocity by
$\lambda = \sqrt{K_0/K}$, with $K_0 = \frac12N_{\mathrm f}\kB T_0$ the
kinetic energy of the target temperature $T_0$; multiplying every
velocity by $\lambda$ multiplies $K$ by $\lambda^2$. The kinetic energy is
then exactly $K_0$ at the start of every step, and the only fluctuation
left is what the motion produces within one step: over \SI{95}{\pico
\second} of liquid argon, recorded at the ends of the steps, the relative
spread of $K$ is 0.0056, against the canonical 0.0511. The canonical
ensemble has a spread; rescaling all but removes it, and with it any use
of the fluctuation formulae of \cref{sec:sm-kinetic}. The scaling factor $\lambda$ is a
symbol of this chapter only (\cref{sec:la-multipliers,sec:cs-shake} used
it for a multiplier).

\subsection*{Weak coupling}

Berendsen and colleagues asked the temperature instead to relax towards
$T_0$ with a time constant $\tau_{\mathrm T}$\cite{Berendsen1984},
\begin{equation*}
   \frac{\dd T}{\dd t} = \frac{T_0 - T}{\tau_{\mathrm T}} .
\end{equation*}
Here $T$ is the kinetic temperature $T_{\mathrm{kin}}$ of Chapter~12.
Over one step $\dt$ the kinetic temperature should change by $(\dt/
\tau_{\mathrm T})(T_0 - T)$. Scaling the velocities by $\lambda$ turns $T$
into $\lambda^2T$, so
\begin{equation}
   \lambda^2T = T + \frac{\dt}{\tau_{\mathrm T}}(T_0 - T) , \qquad
   \lambda^2 = 1 + \frac{\dt}{\tau_{\mathrm T}}\Bigl(\frac{T_0}{T} - 1
   \Bigr) .
   \label{eq:th-berendsen}
\end{equation}
At \SI{150}{\kelvin}, with a target of \SI{135}{\kelvin}, steps of
\SI{10}{\femto\second} and $\tau_{\mathrm T} = \SI{100}{\femto\second}$,
$\lambda^2 = 0.9900$ and $\lambda = 0.99499$: the velocities shrink by half
a percent. With $\tau_{\mathrm T} = \dt$ this is rescaling; as
$\tau_{\mathrm T}$ grows, the run approaches one at fixed energy.

One thing follows from \cref{eq:th-berendsen}. The scaling turns $K$ into
$\lambda^2K$, so each step adds the heat $(\lambda^2 - 1)K =
(\dt/\tau_{\mathrm T})(K_0 - K)$, since $T_0K/T = K_0$. Over a long run in
which nothing drifts this heat averages to zero, because by
\cref{eq:th-effective} the total heat is the change of $K + U$, which stays
within its fluctuations however long the run; so $\ens{K} = K_0$ and the
mean temperature is exactly the target. Over three runs it is \SI{135.00}{\kelvin} with
$\tau_{\mathrm T} = \SI{100}{\femto\second}$ and \SI{135.02}{\kelvin}
with \SI{1}{\pico\second}. But the spread of $K$ is not the canonical
one (\cref{fig:th-berendsen}(a)): with $\tau_{\mathrm T} =
\SI{100}{\femto\second}$ it is 0.0228, 0.45 of the canonical spread and
smaller even than the 0.0311 of a run at fixed energy; with
\SI{1}{\pico\second}, 0.0300, close to the run at fixed energy. Berendsen
coupling does not sample the canonical ensemble at any $\tau_{\mathrm T}$:
strong coupling squeezes the fluctuations, and weak coupling leaves them
those of fixed energy.

\begin{figure}[htb]
\centering
\includegraphics{ch13_thermostats/berendsen}
\caption{(a) The kinetic energy of liquid argon, divided by its mean,
over three runs of \SI{95}{\pico\second} under Berendsen and under CSVR
(\cref{sec:th-csvr}), both with $\tau_{\mathrm T} =
\SI{100}{\femto\second}$, against the canonical gamma density of 765
freedoms (dashed). (b) The liquid with its velocities scaled to
\SI{200}{\kelvin}, brought towards \SI{135}{\kelvin} by Berendsen
and CSVR with $\tau_{\mathrm T} = \SI{1}{\pico\second}$, a Nosé-Hoover
chain with $\tau_{\mathrm T} = \SI{1}{\pico\second}$
(\cref{sec:th-nose}) and Langevin with $\gamma = \SI{1}{\per\pico\second}$
(\cref{sec:th-langevin}), against the approach predicted for Berendsen
(thin). The target is the dotted line.}
\label{fig:th-berendsen}
\end{figure}

\subsection*{What the time constant means}

The temperature takes longer than $\tau_{\mathrm T}$ to settle.
\Cref{fig:th-berendsen}(b) starts the liquid at \SI{200}{\kelvin} by
scaling its velocities. Within \SI{0.2}{\pico\second} the kinetic
temperature falls to between 165 and \SI{170}{\kelvin} under every
thermostat, as kinetic energy passes into potential energy, the sharing
of \cref{sec:sm-equipartition}. After that the thermostat removes heat at
the rate $(K - K_0)/\tau_{\mathrm T} = C_K(T - T_0)/\tau_{\mathrm T}$,
with $C_K = \frac12N_{\mathrm f}\kB$ the kinetic part of the heat
capacity (\cref{sec:sm-kinetic}). But the liquid's energy exceeds its
value at $T_0$ by $C_V(T - T_0)$, more than the $C_K(T - T_0)$ held as
kinetic energy, so it
changes at the rate $C_V\,\dd T/\dd t$, and equating the two,
\begin{equation*}
   C_V\frac{\dd T}{\dd t} = -\frac{C_K}{\tau_{\mathrm T}}(T - T_0) ,
\end{equation*}
and $T - T_0$ decays as $\mathrm{e}^{-t/\tau_{\mathrm{eff}}}$, as the gap
$w$ of \cref{sec:ne-drag} did, with $\tau_{\mathrm{eff}} = \tau_{\mathrm
T}C_V/C_K$. With $C_V = 2.376\,N\kB$ and $C_K = 1.4941\,N\kB$ from
\cref{sec:sm-kinetic}, $C_V/C_K = 1.590$, so $\tau_{\mathrm{eff}} =
\SI{1.59}{\pico\second}$ for $\tau_{\mathrm T} = \SI{1}{\pico\second}$; a
fit to the run from \SI{0.5}{\pico\second} on gives
\SI{1.67}{\pico\second}, uncertain by \SI{0.19}{\pico\second} in a single
run.

\subsection*{The flying ice cube}

Rescaling and Berendsen coupling multiply every velocity by the same
factor, including the part of each velocity that belongs to the motion of
the whole box. No force changes that motion (\cref{sec:ne-atoms}), so its
kinetic energy $K_{\mathrm{com}}$ changes only by the thermostat, by the
factor $\lambda^2$ at each step. Over $n$ steps, $\ln K_{\mathrm{com}}$
gains $\sum\ln\lambda^2$. Writing $T = T_0 + \delta T$, whose departure
$\delta T$ from the target has mean zero since $\ens{T} = T_0$,
\begin{align*}
   \lambda^2 - 1 &= \frac{\dt}{\tau_{\mathrm T}}\Bigl(\frac{1}{1 + \delta
   T/T_0} - 1\Bigr) = \frac{\dt}{\tau_{\mathrm T}}\Bigl(-\frac{\delta T}{T_0}
   + \Bigl(\frac{\delta T}{T_0}\Bigr)^2 - \dotsb\Bigr) ,\\
   \ens{\ln\lambda^2} &\approx \ens{\lambda^2 - 1}
   \approx \frac{\dt}{\tau_{\mathrm T}}\Bigl\langle\Bigl(\frac{\delta
   T}{T_0}\Bigr)^2\Bigr\rangle ,
\end{align*}
to first order in $\dt/\tau_{\mathrm T}$ and second in $\delta T/T_0$,
using $1/(1 + y) = 1 - y + y^2 - \dotsb$, which holds because $(1 + y)(1 -
y + y^2) = 1 + y^3$, and $\ln(1 + y) \approx y$ for small $y$, the Taylor
series of \cref{sec:mo-taylor} with the slope $1/(1 + y)$ of the
logarithm, 1 at $y = 0$. The mean is positive whatever the sign of the
fluctuations, and over $n$ steps, a time $t = n\dt$, $\ln K_{\mathrm{com}}$
gains $n\ens{\ln\lambda^2} \approx \ens{(\delta T/T_0)^2}t/\tau_{\mathrm
T}$: $K_{\mathrm{com}}$ grows as $\mathrm{e}^{\ens{(\delta
T/T_0)^2}t/\tau_{\mathrm T}}$, at
the expense of the motion of the atoms against one another, which the
thermostat cools to keep the total at $K_0$. In time the box flies
through space while its atoms freeze, the flying ice cube of Harvey, Tan
and Cheatham, who found velocity rescaling to break the equipartition of
energy\cite{Harvey1998}, revisited by Braun, Moosavi and
Smit\cite{Braun2018}.
For liquid argon with $\tau_{\mathrm T} = \SI{100}{\femto\second}$,
$\ens{(\delta T/T_0)^2} = \num{4.99e-4}$ in the first of the three runs, and the
predicted rate is \SI{0.0050}{\per\pico\second}; keeping the next term of
$\ln(1 + y)$, $-\frac12y^2$, lowers it by the factor $1 - \dt/2\tau_{\mathrm
T} = 0.95$, to \SI{0.0047}{\per\pico\second}, which
\cref{sec:th-trajectory} measures. A weaker form of the same flow feeds
any motion that the forces barely touch, such as the slow vibrations of
large molecules, held back only by its slow exchange of energy with the
other motions, which is why the effect appears even when the motion of
the centre of mass has been removed\cite{Harvey1998,Braun2018}. For one lithium atom on the surface, whose two freedoms are
its whole motion, Berendsen coupling with $\tau_{\mathrm T} =
\SI{100}{\femto\second}$ drives the atom within a picosecond or two onto
a circular orbit of radius \SI{0.48}{\angstrom} round its hollow, on which
$K$ stays at $\kB T_0$ and $K + U$ at \SI{0.19}{\electronvolt}, short of
the \SI{0.3}{\electronvolt} barrier, and the thermostat has nothing left to
correct: the 1000 atoms make 0.001 hops per atom per picosecond, against
0.93 to 1.09 under the thermostats that sample correctly.

\begin{practice}
Rescaling and Berendsen coupling are for bringing a system near a
temperature quickly. Their averages of $T$ are right, but their
fluctuations are not, and a long run under them feeds slow motions and
the motion of the centre of mass. A production run, whose fluctuations or slow motions matter,
uses one of the thermostats of the next sections.
\end{practice}

\begin{keyidea}
Berendsen coupling scales the velocities by the factor $\lambda$ of
\cref{eq:th-berendsen}, so that $T$ relaxes towards $T_0$.
It holds the mean temperature exactly but narrows the spread of $K$
below the canonical one, and it feeds motions that the forces do not
touch, at the rate $\ens{(\delta T/T_0)^2}/\tau_{\mathrm T}$: the flying
ice cube. The
temperature of a liquid relaxes with the time constant $\tau_{\mathrm T}
C_V/C_K$, longer than $\tau_{\mathrm T}$.
\end{keyidea}

\notebook{13}{set the coupling time and watch the spread of K and the drift}

\begin{selfcheck}
What is $\lambda$ for a step of \SI{2}{\femto\second} with $\tau_{\mathrm
T} = \SI{1}{\pico\second}$, at \SI{310}{\kelvin} with a target of
\SI{300}{\kelvin}?
\end{selfcheck}
